\subsection{The level-zero class lattice and its metric}
\label{subsec:selected-column-metric}

The integral terminal values enter a finite class sum. The source of this sum carries a level-zero metric with class masses, which must be kept when its determinant is taken.

Fix the selected sequence \(u\), write \(T=T_u\), and let \(E\) be the CM field of degree \(2f\) and \(K\) the coefficient field.
At level zero the underlying Hermitian space is positive definite
at every embedding and has rank \(k\) over \(E\). The row spaces
of mixed signature do not enter the following construction.
Use the level given by the word marking and the two type primes
in Section~\ref{sec:ramified-column}. Its class set is nonempty
and its rational stabilizers are trivial by the argument there.

\begin{proposition}[The definite column with its finite tests and Hermitian metric]\label{prop:level-zero-column-coordinates}
For the selected sequence at level zero, each class contributes independent rational coefficients. Its discriminant lattice in \eqref{eq:selected-column-discriminant} satisfies
\[
 L_c^\vee/L_c\simeq{\cal O}_E/T{\cal O}_E,\qquad [L_c^\vee:L_c]=T^{2f}.
\]
 Their change
to \(Q_{c\tau}\) identifies the given integral test with membership
in the divided-root Weyl lattice at each finite place. Under the
same change the squared norm is \eqref{eq:selected-column-gram},
with an orthogonal decomposition over the classes.
\end{proposition}

\begin{proof}
\proofstep{Use the marked discriminant quotient.}

For each class \(c\), write \(L_c\) for its underlying Hermitian
lattice. It is the intersection of the given adelic local lattices with the
rational space. The level carries an integral model. Away from \(T\)
the local lattice is self-dual, and at \(T\) the split lattice in the
two reciprocal factors is given, after isometry, by the standard
lattice in one factor and a row kernel of index \(T\) in the other.
Its Hermitian dual interchanges the reciprocal basis matrices.
In particular
\begin{equation}\label{eq:selected-column-discriminant}
 L_c\subset L_c^\vee,\qquad
 L_c^\vee/L_c\simeq \mathcal O_E/T\mathcal O_E,\qquad
 [L_c^\vee:L_c]=T^{2f}.
\end{equation}
Indeed the paired local quotient has one residue line in each
of the two factors. These are all the local quotients, and isometries
of the adelic model preserve the quotient. This uses the
discriminant module in the word marking.
\label{cor:selected-column-metric-1}

\proofstep{Specify the integral coefficient representation.}

The weight at label \(\tau\) is the signed dominant \(k\)-tuple
\(\lambda_\tau\) given by the recorded column weight. Set
\(a_\tau=\lambda_{\tau,k}\),
\(\nu_{\tau,j}=\lambda_{\tau,j}-a_\tau\) and
\(d_\tau=\sum_j\nu_{\tau,j}\). Use the polynomials on a matrix \(Z\)
of \(k\) rows and \(k\) columns generated by the divided root
operators from the highest polynomial
\[
 h_\nu(Z)=\prod_{j=1}^k\det Z_{[1,j],[1,j]}^{\,\nu_j-\nu_{j+1}},
 \qquad \nu_{k+1}=0.
\]
The integral Weyl lattice \(W_\nu\) is the span of these divided
root words, with their integral coefficients. In its rational
span use the action \(R_\lambda(g)P(Z)=\det(g)^aP(g^t Z)\).
The tensor over \(\tau\) is denoted by \(R(g)\). Thus local
unimodular basis changes preserve the corresponding integral
Weyl tensor, including its determinant characters.

\proofstep{Identify the rational class coordinates.}

In a fixed rational frame let \(G\) be the Hermitian matrix
of the level-zero space. Let \(\rho_\tau:E\to K\) be the conjugate
of the embedding used by the form coefficient at \(\tau\).
Write \(Q_{c\tau}\) for \(L_c\) extended along \(\rho_\tau\),
with its extended metric. Use a \(K\)-frame and let \(B_{c\tau}\)
be the matrix of the scalar-extension map in these coordinates.
If the raw coefficient of the form at \(c\) is \(P_c\), its
reading in the \(Q\)-coordinates is
\[
   F_c=R(g_c)P_c,\qquad
   g_{c\tau}=B_{c\tau}\rho_\tau(G)^{-1}.
\]
The raw coefficients are independent class values. Indeed
evaluation at orbit representatives determines the form by its
transformation rule, and arbitrary values extend by the rule
because the stabilizers are trivial. Invertibility of each
\(R(g_c)\) gives the same assertion for \(F_c\).

\proofstep{Read the dual integral basis in the conjugate embedding.}

At a coefficient place \(w\), take a local basis \(Y_c\) of
\(L_c\) at the primes above its residue characteristic, using
one semilocal basis if several primes of \(E\) are involved.
In the \(Q\)-coordinates the basis is

\[
  \gamma_{c\tau}=B_{c\tau}\rho_\tau(Y_c).
\]
 In the raw form
coordinates the dual lattice has basis \((Y_c^*G)^{-1}\), where
\(*\) denotes conjugate transpose over \(E\). Since \(\rho_\tau\)
is the conjugate of the form embedding,
\((\gamma_{c\tau}^{-1}g_{c\tau})^t\) in the coefficient reading
is the form embedding of \((Y_c^*G)^{-1}\). In fact

\[
  \gamma_{c\tau}^{-1}g_{c\tau}=\rho_\tau(GY_c)^{-1},
\]
 and
\(G=G^*\). Transposition preserves the determinant character.
Hence the integral raw form test in the dual basis is precisely
\[
  R(\gamma_c^{-1})F_c\in
       \bigotimes_\tau\bigl(\mathcal O_{K,w}\otimes_\mathbb Z
                              W_{\nu_\tau}\bigr)
  \quad\text{for every }c.
\]
Testing another local basis gives the same result by unimodular
covariance. Intersecting the local tests recovers the full
level-zero lattice \(M_u\).

\proofstep{Transport the squared norm by the matching Gram.}

At an infinite place \(s\) write \(\|v\|_{Q_{c\tau},s}
=\|U_{c\tau,s}v\|_2\) in the chosen \(Q\)-frame. Thus
\[
  (U_{c\tau,s}B_{c\tau}^{\,s})^*
     (U_{c\tau,s}B_{c\tau}^{\,s})=s\rho_\tau(G).
\]
It follows that
\((U_{c\tau,s}g_{c\tau}^{\,s})^t\) is an orthonormal frame for
the raw form Gram \(s\rho_\tau(G)^t\). This is the same
conjugate-transpose relation as at the finite place. The
class-weight in the level-zero norm is one, since the stabilizer
is trivial and the level-zero measure weight is one.

The use of
an orthonormal frame separately for each \(c\) is valid because
one can test the norm identity on forms supported at \(c\).
Equivalently the normalized Fock form is invariant under unitary
changes of the frame, including the absolute value of the
determinant character. Thus the squared norm is
\begin{equation}\label{eq:selected-column-gram}
   \|x\|_{M_u,s}^2=
   \sum_c\frac{\|R(U_{c,s})F_c^{\,s}\|_{\mathrm F}^2}
                 {\prod_\tau\|h_{\nu_\tau}\|_{\mathrm F}^2},
   \qquad
   \|P\|_{\mathrm F}^2=\sum_\alpha\alpha!|P_\alpha|^2 .
\end{equation}
The equality of the norms gives the Hermitian pairing by
polarization. In particular the class decomposition is
orthogonal at every embedding. This is the column metric used
for the trace-zero summand and for the later normalized
Gram matrices. \label{cor:selected-column-metric-2}

\end{proof}

\begin{proposition}[Both slope estimates for the level-zero column]\label{prop:level-zero-column-slopes}
Use the signed dominant weights of
Proposition~\ref{prop:level-zero-column-coordinates} and the
discriminant module \eqref{eq:selected-column-discriminant}.
Put \(r_E=|\operatorname{disc}(E)|^{1/(2f)}\) and
\(c_\nu=\|h_\nu\|_{\rm F}^2/\prod_i\nu_i!\) for the
unnormalized Fischer norm. The level-zero lattice satisfies
\[
\begin{aligned}
 \mu(M_u)&\ge -\sum_\tau\frac{\sum_i\lambda_{\tau,i}}{k}\log T
       -\sum_\tau\bigl(d_\tau\log k+\log c_{\nu_\tau}\bigr),\\
 \mu_{\max}(M_u)&\le
    \sum_\tau\bigl(d_\tau\log(k^4r_E^2)+
         \log c_{\nu_\tau}-a_\tau\log T\bigr).
\end{aligned}
\]
These estimates use the same metric and integral value lattices
at all classes.
\end{proposition}

\begin{proof}
In a rational frame, if \(J\) is the determinant ideal of a
Hermitian lattice with squared-norm Gram \(H_s\), the degree
in the convention of the comparison is
\(-2\log\mathrm N J/[K:\mathbb Q]-\operatorname{avg}_s\log\det H_s\).
An isometric base change preserves both normalized terms.
Taking the Hermitian dual inverts this degree, by dualizing the
finite determinant and the real Gram in the same frame and
using the product formula. Consequently
\eqref{eq:selected-column-discriminant} gives
\(\mu(Q_{c\tau})=-\log T/k\). The factor of two in the
finite degree cancels when comparing \(L_c\) and its dual.

\proofstep{Use the determinant character to compute the contribution to the mean slope.}
For a representation of dimension \(r\), the character
\(\det R_\lambda(g)\) is \((\det g)^{r\sum_i\lambda_i/k}\).
If \(b\) is the common exponent on the diagonal coordinates, evaluating at \(tI_k\) gives
\[
  t^{kb}=\det R_\lambda(tI_k)=t^{r\sum_i\lambda_i},\qquad b=\frac r k\sum_i\lambda_i.
\]
One can see the exponent on the diagonal torus, since conjugation
by permutation matrices gives the same exponent on each coordinate,
and a scalar matrix acts by its \(\sum_i\lambda_i\)-th power.

The formula applies both to the finite frames and to the
archimedean changes of frame. Tensoring with the other labels
repeats each matrix according to their dimensions. The metric
on the standard Weyl tensor is the tensor product of the
normalized Fock forms. For a weight \(\nu\), let
\(\mu_{\rm st}(\nu)\) be minus the logarithm of the Gram
determinant of \(W_\nu\), divided by its rank, with \(h_\nu\)
of norm one. Taking the determinant and dividing
the direct sum over \(c\) by its rank gives
\begin{equation}\label{eq:selected-column-mean-exact}
  \mu(M_u)=
     -\sum_\tau\frac{\sum_i\lambda_{\tau,i}}{k}\log T
       +\sum_\tau\mu_{\rm st}(\nu_\tau),
\end{equation}

For the standard term put \(c_\nu=\|h_\nu\|_{\mathrm F}^2/\prod_j\nu_j!\).
The Fock highest-vector formula gives
\[
 c_\nu=\prod_{i<j}\frac{\nu_i+j-i}{\nu_i-\nu_j+j-i},
 \qquad 1\le c_\nu\le(\nu_1+k)^{k(k-1)/2}.
\]
For example, the formula follows by integrating over a square
standard complex Gaussian matrix \(Z=UR\), with \(R\) upper
triangular and positive on its diagonal. Orthogonal projection
of the independent Gaussian columns gives independent
\(R_{jj}^2\) of gamma shape \(k-j+1\), independent of the Haar
unitary \(U\). Each leading minor of \(Z\) is the corresponding
minor of \(U\) times \(\prod_{i\le j}R_{ii}\). The Haar mean of
\(|h_\nu(U)|^2\) is the reciprocal Weyl dimension, by Schur
orthogonality of the unit highest-weight matrix coefficient.
In particular the unitary integral used in the moment calculation is
\[
  \int_{U(k)}|h_\nu(U)|^2\,dU=(\dim V_\nu)^{-1}.
\]
The gamma moments and the Weyl dimension formula give the displayed
product after dividing by \(\prod_j\nu_j!\).

The contravariant integral Weyl duality identifies the Gram
dual of \(W_\nu\) in the normalized Fock form with the
dual Weyl lattice of integral polynomials of column degrees
\(\nu_j\) invariant under the upper unipotent column group.
The normalization here pairs the two highest polynomials to
one. This is the integral duality for the divided-root Weyl
module, realized by the Schur polynomial lattice, and is used
only for the standard coefficient lattice. If a monomial
has these column degrees, 
\[
 \frac{\prod_j\nu_j!}{\alpha!}=\prod_j\binom{\nu_j}{\alpha_{1j},\ldots,\alpha_{kj}}\le\prod_j k^{\nu_j}=k^{\sum_j\nu_j}
\]
 by the multinomial theorem in each column.
A full-rank integer coefficient matrix has a nonzero maximal
minor of absolute value at least one. Cauchy--Binet for the
dual Gram therefore bounds its determinant below by
\((c_\nu k^{\sum_j\nu_j})^{-\operatorname{rk}W_\nu}\).
Inverting the determinant gives
\[
  \mu_{\rm st}(\nu)\ge -\Big(\sum_j\nu_j\Big)\log k-\log c_\nu .
\]
These are the standard slope and the hook estimate in the
lower slope proof. \label{cor:selected-column-metric-3}

\proofstep{Compare to a short-dual frame without changing the local lattice.}

Use the positive real trace form \(b(x,y)=\mathrm{Tr}_{E/\mathbb Q}
h(x,y)\) on the Minkowski space of the Hermitian lattice \(L_c\).
Write \(\mathfrak d\) for the different, \(D_E=|\mathrm{disc}(E)|\)
and \(r_E=D_E^{1/(2f)}=2v\), where \(v\) is the host
discriminant parameter. The \(b\)-dual of \(L_c^\vee\) is
\(\mathfrak d^{-1}L_c\), since \(L_c^\vee\) is an
\(\mathcal O_E\)-module and the double Hermitian dual is \(L_c\).
For any nonzero \(x\) in this trace dual,

\[
  h(x,x)\mathfrak d^2\subset\mathcal O_E,
\]
 using
\(L_c\subset L_c^\vee\) and conjugation invariance of
\(\mathfrak d\). Its field norm is at least \(D_E^{-2}\) and
its values at all embeddings are positive. The arithmetic-geometric
mean inequality gives \(b(x,x)\ge2f r_E^{-2}\).

Apply Banaszczyk's Euclidean successive-minima inequality

\[
  \lambda_n(\Lambda)\lambda_1(\Lambda^*)\le n
\]
 with
\(n=2fk\), \(\Lambda=L_c^\vee\) in the trace norm and
\(\Lambda^*\) its trace dual. Among \(2fk\) real-independent
short vectors choose \(k\) that are \(E\)-independent, since
the full list spans the ambient rational space. We obtain
\(y_1,\ldots,y_k\in L_c^\vee\) with
\[
       (2f)^{-1}\mathrm{Tr}_{E/\mathbb Q}h(y_i,y_i)
                   \le k^2 r_E^2 .
\]
This bound involves no power of \(T\). The integrality
in \eqref{eq:selected-column-discriminant} supplies the
shortest-vector bound on the dual of \(L_c^\vee\).
\cite{banaszczyk}\label{cor:selected-column-metric-4}

Let \(C_{c\tau}\) have rows the Hermitian pairing with the
images of the \(y_i\) in \(Q_{c\tau}\).
These are \(K\)-linear functionals and they form a basis
of the rational dual. For each class read the coefficients of
\(F_c(C_c^tZ)\), with the substitution made on each label.
This is an injective rational map on the class summand.
For a local basis \(\gamma_{c\tau}\) of \(Q_{c\tau}\),
\(C_{c\tau}\gamma_{c\tau}\) is integral.

Write
\(\det Q_{c\tau}\) for the finite determinant ideal in the
chosen frame, with local generator \(\det\gamma_{c\tau}\).
By the finite
coordinate formula, each coefficient of the image belongs to
the fractional ideal
\[
             I_c=\prod_\tau(\det Q_{c\tau})^{a_\tau}.
\]
Indeed substitution by the integral
matrix \(C_{c\tau}\gamma_{c\tau}\) acts on polynomial
coefficients, and the omitted determinant character of
\(C_{c\tau}\) leaves the finite multiplier
\((\det\gamma_{c\tau})^{a_\tau}\).

Give the monomial coefficient module in this ideal the
diagonal squared norm \(\alpha!\) on monomial \(\alpha\),
and divide the squared norm by the common squared
highest-vector norm.
At \(s\) the substitution matrix in the normalized frame and its scalar determinant factor are respectively
\[
  A=(C_{c\tau}^sU_{c\tau,s}^{-1})^t,\qquad b=(\det U_{c\tau,s})^{-a_\tau}.
\]
For the polynomial of degree \(d_\tau\) a substitution has
squared operator norm bounded by the \(d_\tau\)-th power of
the squared Hilbert--Schmidt norm of its matrix. This follows
by realizing the homogeneous Fock space as symmetric tensors
with the induced norm times \(d_\tau!\).

For a substitution matrix \(A\), this comparison reads
\[
  \|\operatorname{Sym}^{d_\tau} A\|_{\rm op}^2\le\|A\|_{\rm op}^{2d_\tau}\le\|A\|_{\rm HS}^{2d_\tau}.
\]
 The squared
Hilbert--Schmidt norm of this matrix is
\[
 h_{c\tau,s}=\sum_i s\rho_\tau\bigl(h(y_i,y_i)\bigr).
\]
The values are positive real numbers. The mean of
\(h_{c\tau,s}\) over the infinite places of \(K\) is the
mean of the analogous sum over \(E\), because both fields
are totally complex and each base embedding has the same
number of extensions. Thus
\(\operatorname{avg}_s\log h_{c\tau,s}\le\log(k^3r_E^2)\).
The remaining height of the map, including the target ideal, is
\[
 -\frac{2\log\mathrm N I_c}{[K:\mathbb Q]}
    -\sum_\tau a_\tau\operatorname{avg}_s\log|\det U_{c\tau,s}|^2
 =\sum_\tau a_\tau\deg Q_{c\tau}
 =-\sum_\tau a_\tau\log T .
\]
The frame metric and determinant ideal are used together in this
identity, including for negative \(a_\tau\).

The lower bound on \(\alpha!/\prod_\tau
\|h_{\nu_\tau}\|_{\mathrm F}^2\) from the column degrees
bounds the maximum slope of the diagonal coefficient module.
For clarity, an integral lattice in an orthogonal direct sum of
Hermitian lines has maximum slope at most the largest line
slope in that sum. This follows by projecting successively to the
lines and using the quotient norms, or by taking determinants
of the projections in a filtration. A \(K\)-linear injection
into that sum bounds every source slope above by the largest
target slope plus the averaged logarithm of the squared
operator bound.

For a rank-\(r\) sublattice the norm estimate on its determinant follows from
\[
  \|\textstyle\bigwedge^r A\|_{\rm op}\le\|A\|_{\rm op}^r .
\]
 This follows on every sublattice by applying
the bound to its determinant, with the finite index of its
image included in the target degree. Combining the last three
bounds gives, on each class,
\[
  \mu_{\max}(M_{u,c})\le
    \sum_\tau\bigl[d_\tau\log(k^4r_E^2)
                       +\log c_{\nu_\tau}-a_\tau\log T\bigr].
\]
The same bound holds on the orthogonal integral sum. This
reasoning bounds the logarithmic height over all infinite
places, so it allows unequal pointwise lengths at the
embeddings. It is the short-dual comparison on the actual
column Gram. \label{cor:selected-column-metric-5}

\proofstep{Insert the signed top and non-top weight bounds.}
To apply the two estimates, at a top label the first entry is
in \([a,2a]\), the middle entries lie in \([0,S_c]\), and the
last entry is in \((-2b',-b']\). At the other labels all
entries lie in \([0,C_n S_c]\) for a fixed size constant \(C_n\).
It follows that

\[
  \sum_\tau d_\tau\le 2|T_{\rm top}|a+
O_{C_n}(fk(b'+S_c)),
\]
 and
\(\sum_\tau|a_\tau|=O_{C_n}(f(b'+S_c))\).
The sum of the signed weights is at most
\(\sum_{t\in T_{\rm top}}a_t(u)+O_{C_n}(fkS_c)\).
On the size range, \(\log(\nu_{\tau,1}+k)\ll L\) and
\(\log(k^4r_E^2)\le4\log(kv)\), the latter since \(v\ge2\).
Substituting these estimates gives the two inequalities of
Section~\ref{sec:frames}. Their constants are uniform over
classes, words, and coefficient places once the common size
constants are fixed. The \(T_{\rm top}\) in these size
formulas is the set of top labels, whereas \(T=T_u\) in the
logarithms is the prime assigned to the sequence.
\end{proof}

